%!TEX program = xelatex
\documentclass[11pt,a4paper]{article}
\usepackage[margin=2.2cm]{geometry}
\usepackage{amsmath,amssymb,amsthm,mathtools}
\usepackage{bm}
\usepackage{enumitem}
\usepackage{booktabs}
\usepackage{xcolor}
\usepackage{hyperref}
\usepackage{fancyhdr}
\usepackage{xeCJK}

\hypersetup{
    colorlinks=true,
    linkcolor=blue!60!black,
    citecolor=blue!60!black,
    urlcolor=blue!60!black
}

\pagestyle{fancy}
\fancyhf{}
\rhead{\small Qiuzhen QE Computational \& Applied Math --- Fall 2023}
\lhead{\small Official Qualifying Exam}
\rfoot{\small Page \thepage}

\DeclareMathOperator*{\argmax}{arg\,max}
\DeclareMathOperator*{\argmin}{arg\,min}

\setlist[itemize]{topsep=3pt,itemsep=2pt,parsep=0pt}
\setlist[enumerate]{topsep=3pt,itemsep=2pt,parsep=0pt}

\title{\textbf{QUALIFYING EXAM: APPLIED MATHEMATICS}\\[0.5em]\Large SEPTEMBER 2023}
\date{}
\author{}

\begin{document}

\maketitle
\vspace{-3em}

\noindent\rule{\linewidth}{0.4pt}

\begin{enumerate}[label=\textbf{\arabic*.}]
    \item Consider Newton's method for finding a solution $x_*$ to $f(x) = 0$, where $f \in C^2(a, b), x_* \in (a, b)$.
    \begin{align*}
    1: & \quad \text{Determine } x_0 \in (a, b) \\
    2: & \quad \text{for } k = 0, 1, 2, \dots \text{ do} \\
    3: & \qquad x_{k+1} = x_k - \frac{f(x_k)}{f'(x_k)} \\
    4: & \quad \text{end for}
    \end{align*}
    \begin{enumerate}[label=(\arabic*)]
        \item Prove that if $x_0$ is sufficiently close to $x_*$ and $f'(x_*) \neq 0$, then $\lim_{k \to \infty} x_k = x_*$ and
        \[
        \lim_{k \to \infty} \frac{x_{k+1} - x_*}{(x_k - x_*)^2} = \frac{f''(x_*)}{2 f'(x_*)}.
        \]
        \item In practice, the derivative is sometimes not easy to obtain. As a result, secant differences are used instead:
        \[
        x_{k+1} = x_k - \frac{x_k - x_{k-1}}{f(x_k) - f(x_{k-1})} f(x_k).
        \]
        Prove that if $x_0, x_1$ are sufficiently close to $x_*$, then $x_k \to x_*$ and
        \[
        \lim_{k \to \infty} \frac{x_{k+1} - x_*}{(x_k - x_*)(x_{k-1} - x_*)} = \frac{f''(x_*)}{2 f'(x_*)}.
        \]
    \end{enumerate}

    \item Consider the explicit shifted QR method for computing eigenvalues of $A \in \mathbb{C}^{n \times n}$:
    \begin{align*}
    1: & \quad \text{Find upper Hessenberg } H_0 \text{ and unitary } U_0 \text{ such that } H_0 = U_0^H A U_0 \\
    2: & \quad \text{for } k = 0, 1, 2, \dots \text{ do} \\
    3: & \qquad \text{Determine a scalar } \mu_k \\
    4: & \qquad \text{Compute QR factorization } Q_k R_k = H_k - \mu_k I \\
    5: & \qquad H_{k+1} = R_k Q_k + \mu_k I \\
    6: & \quad \text{end for}
    \end{align*}
    \begin{enumerate}[label=(\arabic*)]
        \item Prove that $H_i, i = 1, 2, \dots$ are all upper Hessenberg matrices.
        \item Interpret the purpose of using $H_0$ rather than $A$ in the iteration.
        \item Suppose $A$ has $n$ distinct eigenvalues and none of the shifts $\mu_i$ is an eigenvalue of $A$. Prove that $H_i$ are unreduced upper Hessenberg matrices.
        \item Write $H_k = \begin{bmatrix} G_k & u_k \\ \varepsilon_k e^T & \alpha_k \end{bmatrix}$ where $\alpha_k, \varepsilon_k \in \mathbb{C}$, $u_k, e \in \mathbb{C}^{n-1}$ with $e = [0, \dots, 0, 1]^T$. Prove that
        \[
        |\varepsilon_{k+1}| \le \rho_k^2 \|u_k\|_2 |\varepsilon_k|^2 + \rho_k |\alpha_k - \mu_k| |\varepsilon_k|,
        \]
        where $\rho_k = \|(G_k - \mu_k I)^{-1}\|_2$, provided $\mu_k$ is not an eigenvalue of $G_k$.
    \end{enumerate}

    \item If $p$ is a polynomial of degree $n$ on $[-1, 1]$, it is determined by its values on an $(n+1)$-point grid on $[-1, 1]$. The derivative $p'$, a polynomial of degree $n-1$, is determined on the same grid. The differentiation matrix $D = (D_{ij}) \in \mathbb{R}^{(n+1) \times (n+1)}$ represents the linear map $p'(x_i) = \sum_{j=0}^n D_{ij} p(x_j)$.
    \begin{enumerate}[label=(\arabic*)]
        \item Prove that $D_{ij} = l_j'(x_i)$, where $l_j(x)$ is the $j$-th Lagrange basis function.
        \item If a Chebyshev grid $x_j = \cos(j\pi/n), 0 \le j \le n$ is adopted, derive explicit formulas for $D_{ij}$.
    \end{enumerate}

    \item Consider a Runge-Kutta method for $y' = f(t, y)$ with Butcher tableau:
    \[
    \begin{array}{c|cc}
    0 & 0 & 0 \\
    1 & 1 & 0 \\
    \hline
    & 1/2 & 1/2
    \end{array}
    \]
    \begin{enumerate}[label=(\arabic*)]
        \item Rewrite the scheme in the form $u_{n+1} = u_n + h F(t_n, u_n, h; f)$.
        \item Assume $f$ is sufficiently smooth. Prove that this method is convergent and determine the order of convergence.
        \item What is the region of absolute stability of this method? Give an explicit description.
    \end{enumerate}

    \item For $u_t + a u_{xxx} = 0$ ($a$ constant), applying Lax-Friedrichs idea yields:
    \[
    u_m^{n+1} = \frac{1}{2}(u_{m+1}^n + u_{m-1}^n) - \frac{1}{2} a k h^{-3} (u_{m+2}^n - 2u_{m+1}^n + 2u_{m-1}^n - u_{m-2}^n).
    \]
    \begin{enumerate}[label=(\arabic*)]
        \item Give the leading order term of the local truncation error.
        \item Analyze the stability of this scheme.
    \end{enumerate}

    \item Consider linear programming $\max_{\mathbf{x}} \langle \mathbf{c}, \mathbf{x} \rangle$ s.t. $A\mathbf{x} = \mathbf{b}, \mathbf{x} \ge 0$, where $\mathbf{b} \ge 0$, $A \in \mathbb{R}^{m \times n}$ has full rank ($m < n$).
    \begin{enumerate}[label=(\arabic*)]
        \item Prove that basic feasible solutions are equivalent to the vertices of its feasible region.
        \item Write down the dual problem, and prove strong duality when an optimal solution exists.
    \end{enumerate}

    \item Consider the eigenvalue problem with $0 < \epsilon \ll 1$:
    \[
    u'' + (\lambda + \epsilon f(x))u = 0, \quad 0 < x < 1; \qquad u(0) = 0, \quad u'(1) = 0.
    \]
    Give the asymptotic expansion of $\lambda$ up to $O(\epsilon)$.
\end{enumerate}

\end{document}
